% Part: first-order-logic
% Chapter: model-theory
% Section: nonstandard-arithmetic

\documentclass[../../../include/open-logic-section]{subfiles}

\begin{document}

\olfileid{mod}{bas}{nsa}
\section{Modèl Aritmetika Nonstandar}

\begin{defn}
Anggep $\Lang{L_N}$ minangka !!{language} aritmetika,
kang ngemot !!{constant} $\Obj{0}$, !!{predicate} rong panggonan
$<$, !!{function} siji panggonan $\prime$, lan
!!{function}s rong panggonan $+$ lan $\times$.
\begin{enumerate}
\item \emph{Modèl standar} aritmetika yaiku !!{structure}
  $\Struct{N}$ kanggo $\Lang{L_N}$ kanthi
  $\Nat = \{ 0, 1, 2, \dots\}$, kang napsiraké
  $\Obj{0}$ dadi $0$, $<$ dadi relasi kurang saka
  ing~$\Nat$, lan $\prime$, $+$, lan $\times$
  dadi fungsi penerus, panjumlahan, lan perkalian
  ing $\Nat$ miturut urutané.
\item \emph{Aritmetika bener} yaiku téori $\Theory{N}$.
\end{enumerate}
\end{defn}

Nalika makarya ing $\Lang{L_N}$, saben suku awujud
$\Obj{0}^{\prime\dots\prime}$, kanthi $n$ panganggoné
fungsi penerus marang $\Obj{0}$, dicekak dadi~$\num{n}$.

\begin{defn}
!!{structure}~$\Struct{M}$ kanggo $\Lang{L_N}$ diarani
\emph{standar} yèn lan mung yèn $\Struct{N} \iso \Struct{M}$.
\end{defn}

\begin{thm}\ollabel{thm:non-std}
Ana modèl nonstandar kang !!{enumerable} kanggo aritmetika bener.
\end{thm}

\begin{proof}
Ambakna $\Lang{L_N}$ kanthi nambah !!{constant} anyar~$c$,
banjur delengen téori
\[
\Theory{N} \cup \Setabs{\num{n} < c}{n \in \Nat}.
\]
Téori iki bisa dicukupi kanthi winates, mula miturut
kompaknes nduwèni modèl $\Struct{M}$.
Miturut téoréma L\"owenheim--Skolem mudhun, modèl iki bisa
dipilih kang !!{enumerable}. Anggep $\Domain{M}$ domainé
$\Struct{M}$. Ing $\Struct{M}$, lambang nonlogis saka
$\Lang{L_N}$ ditafsiraké dadi
$\mathbf{z} = \Assign{\Obj{0}}{M} \in \Domain{M}$,
${\prec} = \Assign{<}{M} \subseteq \Domain{M}^2$,
$* = \Assign{\prime}{M} \colon \Domain{M} \to \Domain{M}$,
lan $\oplus = \Assign{+}{M}, \otimes =
\Assign{\times}{M} \colon \Domain{M}^2 \to \Domain{M}$.
Kanggo saben $x \in \Domain{M}$, tulis $x^*$ kanggo
unsur $\Domain{M}$ kang diasilaké saka $x$ kanthi
ngetrapaké~$*$.

Saiki, yèn $h$ isomorfisme antarané $\Struct{N}$ lan
$\Struct{M}$, mesthi ana $n \in \Nat$ saéngga
$h(n) = \Assign{c}{M}$. Anggep $s$ sawijining pangaji
ing $\Struct{N}$ kang nyukupi $s(x) = n$. Mula
$\Sat{N}{\eq[\num{n}][x]}[s]$; miturut bukti
\olref[iso]{thm:isom}, uga
$\Sat{M}{\eq[\num{n}][x]}[h \circ s]$, saéngga
$\Assign{c}{M} = \mathbf{z}^{*\cdots *}$
(kanthi $*$ diulang $n$ kaping). Nanging iki ora mungkin,
amarga miturut anggepan $\Sat{M}{\num{n} < c}$
lan $\prec$ iku irrefleksif. Mula $\Struct{M}$ nonstandar.
\end{proof}

\begin{prob}
Relasi $R$ ing himpunan $X$ diarani \emph{berasas apik} yèn lan
namung yèn ora ana ranté mudhun tanpa wates ing~$R$,
yaiku ora ana $x_0$, $x_1$, $x_2$,~\dots ing $X$ saéngga
$\dots x_2Rx_1Rx_0$. Kanthi anggepan téori himpunan
Zermelo--Fraenkel~$ZF$ konsisten, tuduhna yèn ana
modèl $ZF$ kang saka njabané ora berasas apik,
yaiku modèl $\mathfrak{M}$ kang nduwèni
$\dots x_2 \in x_1 \in x_0$.
\end{prob}

Amarga modèl nonstandar $\Struct{M}$ saka \olref{thm:non-std}
ekuivalen sacara elementer karo modèl standar,
sawetara sipat $\Struct{M}$ bisa diturunké.
Pérangan sabanjuré ngrembug sipat-sipat iki kanggo
ngolèhaké gambaran kang cetha bab modèl nonstandar
$\Theory{N}$ kang !!{enumerable}.

\begin{enumerate}
\item Ora ana unsur $\Domain{M}$ kang $\prec$-kurang saka
  awaké dhéwé: ukara $\lforall[x][\lnot x < x]$ bener
  ing $\Struct{N}$ lan mula uga ing $\Struct{M}$.
\item Kanthi alesan kang padha, $\prec$ dadi
  \emph{urutan linier} ing $\Domain{M}$,
  yaiku relasi total, irrefleksif, lan transitif
  ing $\Domain{M}$.
\item Unsur $\mathbf{z}$ iku unsur paling cilik
  miturut $\prec$ ing $\Domain{M}$.
\item Saben unsur $\Domain{M}$ luwih cilik miturut
  $\prec$ tinimbang penerus-$*$-é, lan $x^*$ iku
  unsur paling cilik miturut $\prec$ ing $\Domain{M}$ kang luwih gedhé
  tinimbang $x$.
\item $\Struct{M}$ ngemot pérangan wiwitan miturut
  $\prec$ kang isomorfik karo $\Nat$:
  $\mathbf{z}, \mathbf{z}^*, \mathbf{z}^{**}, \dots$.
  Iki diarani \emph{pérangan standar} saka $\Domain{M}$.
  Saben unsur liyane ing $\Domain{M}$ diarani
  \emph{nonstandar}. Mesthi ana unsur nonstandar ing
  $\Domain{M}$, amarga yèn ora, fungsi $h$ saka
  bukti \olref{thm:non-std} bakal dadi isomorfisme.
  Kita migunakaké $n, m, \dots$ minangka !!{variable}s
  kang ngliwati pérangan standar $\Struct{M}$ iki.
\item Saben unsur nonstandar luwih gedhé tinimbang
  unsur standar endi waé; amarga kanggo saben $n \in \Nat$,
  \[
  \Sat{N}{\lforall[z][(\lnot(\eq[z][\Obj{0}] \lor
  \dots \lor \eq[z][\num{n}]) \lif \num{n} < z)]},
  \]
  mula yèn $z \in \Domain{M}$ béda saka kabèh unsur
  standar, unsur iki mesthi \emph{luwih gedhé} tinimbang
  kabèh unsur standar mau.
\item Saben unsur $\Domain{M}$ saliyané $\mathbf{z}$
  iku penerus-$*$ saka siji unsur unik ing $\Domain{M}$,
  kang dilambangaké $^*x$. Yèn $x = y^*$, mula
  $x$ lan $y$ padha-padha standar yèn salah sijiné
  standar (lan padha-padha nonstandar yèn salah sijiné
  nonstandar).
\item Tetepna relasi ekuivalènsi $\approx$ ing
  $\Domain{M}$: $x \approx y$ yèn lan mung yèn
  ana $n$ \emph{standar} saéngga
  $x \oplus n = y$ utawa $y \oplus n =x$.
  Tegesé, $x \approx y$ yèn lan mung yèn
  jarak antarané $x$ lan $y$ winates.
  Yèn $n$ lan $m$ standar, mula $n \approx m$.
  \emph{Blok} saka $x$ yaiku kelas ekuivalènsiné
  $[x] = \Setabs{y}{x \approx y}$.
\item Anggepen $x \prec y$ lan $x \not\approx y$.
  Amarga $\Sat{N}{\lforall[x][\lforall[y][(x < y \lif (x' < y \lor x' =
      y))]]}$, mula $x^* \prec y$ utawa $x^* = y$.
  Pilihan kapindho ora mungkin amarga ndadèkaké
  $x \approx y$, mula $x^* \prec y$.
  Kanthi alesan kang padha, yèn $x \prec y$ lan
  $x \not\approx y$, mula $x \prec {^*y}$.
  Dadi yèn $x \prec y$ lan $x \not\approx y$,
  saben $w \approx x$ luwih cilik miturut $\prec$
  tinimbang saben $v \approx y$. Mula saben blok
  nonstandar $[x]$ mbentuk ranté tanpa wates ing
  loro arahé:
  \[
  \cdots \prec  {^{**}x} \prec {^*}x \prec x \prec x^* \prec x^{**}
  \prec \cdots
  \]
  Ranté iki diarani ranté-$Z$ amarga jinis urutané
  padha karo wilangan bulat.
\item Urutan $\prec$ bisa ditrapaké marang
  blok-blok kang béda: yèn $x \prec y$ lan
  kaloroné ana ing blok kang béda, mula blok $x$ luwih cilik
  tinimbang blok $y$. Blok diarani
  \emph{nonstandar} yèn ngemot unsur nonstandar.
  Blok standar iku blok paling cilik.
\item Ora ana blok nonstandar paling cilik: yèn $y$
  nonstandar, ana $x \prec y$ kanthi $x$ uga nonstandar
  lan $x \not\approx y$. Buktiné: ing modèl standar
  $\Struct{N}$, saben wilangan bisa dipérang loro,
  bisa uga kanthi sisa siji, yaiku
  $\Sat{N}{\lforall[y][\lforall[x][(y = x + x \lor y = x + x +
        \Obj{0}')]]}$. By elementary equivalence, for every $y \in
  \Domain{M}$ ana $x \in \Domain{M}$ saéngga
  $x \oplus x = y$ utawa
  $x \oplus x \oplus \mathbf{z}^*= y$.
  Yèn $x$ standar, $y$ uga standar; mula $x$
  nonstandar. Kajaba iku, $x$ lan $y$ mapan ing blok
  kang béda, yaiku $x \not\approx y$.
  Kanggo mangertèni iki, anggepen kaloroné ing
  blok kang padha, yaiku $x \oplus n = y$
  kanggo sawatara $n$ standar.
  Yèn $y = x \oplus x$, mula $x \oplus n = x
  \oplus x$, saéngga $x = n$ miturut hukum
  pambatalan kanggo panjumlahan
  (kang bener ing $\Struct{N}$ lan mula uga ing
  $\Struct{M}$). Iki ndadèkaké $x$ standar.
  Kahanan $y = x \oplus x \oplus \mathbf{z}^*$
  padha carané.
\item Kanthi alesan kang padha, ora ana blok paling gedhé.
\item Urutan blok-blok iku rapet: yèn $[x]$ luwih
  cilik tinimbang $[y]$ (kanthi $x \not\approx y$),
  ana blok $[z]$ kang béda saka kaloroné lan mapan
  ing antarané. Anggepen $x \prec y$.
  Kaya mau, $x \oplus y$ bisa dipérang loro
  (bisa uga kanthi sisa), mula ana $u
  \in \Domain{M}$ saéngga
  $x \oplus y = u \oplus u$ utawa
  $x \oplus y = u \oplus u \oplus \mathbf{z}^*$.
  Unsur $u$ iku rata-rata saka $x$ lan $y$, mula
  mapan ing antarané. Anggepen $x \oplus y =
  u \oplus u$ (kahanan sijiné padha carané):
  yèn $u \approx x$, kanggo sawatara $n$ standar,
  \[
  x \oplus y = x \oplus n \oplus x \oplus n,
  \]
  mula $y = x \oplus n \oplus n$ lan kita bakal
  olèh $x \approx y$, kang mbantah anggepan.
  Dadi $u \not\approx x$. Argumèn kang padha
  ngasilaké $u \not\approx y$.
\end{enumerate}
Mula blok-blok nonstandar urutané kaya wilangan rasional:
mbentuk urutan linier kang !!a{enumerable}, rapet,
lan tanpa pucuk. Akibaté, kanggo saben rong modèl
nonstandar !!{enumerable} $\Struct{M}_1$ lan
$\Struct{M_2}$ saka aritmetika bener,
redukté menyang basa kang mung ngemot $<$ lan $=$
padha isomorfik. Isomorfisme $h$ bisa dibangun mangkéné:
pérangan standar $\Struct{M_1}$ lan $\Struct{M_2}$
padha isomorfik karo modèl standar $\Struct{N}$,
mula uga karo siji lan sijiné. Blok-blok kang mbentuk
pérangan nonstandar uga urutané kaya wilangan rasional,
mula miturut \olref[dlo]{thm:cantorQ} padha isomorfik.
Isomorfisme blok iki bisa ditambahi dadi isomorfisme
\emph{ing sajroning} saben blok kanthi nyocogaké
unsur sembarang ing saben pasangan blok, banjur
nyocogaké penerus saka $x$ ing $\Struct{M_1}$
karo penerus saka gambaré $x$ ing $\Struct{M_2}$.
Nanging iki \emph{ora} ateges $\mathfrak{M}_1$
lan $\mathfrak{M}_2$ isomorfik ing basa aritmetika
jangkep (isomorfisme tansah gumantung marang tandha
basa), amarga panjumlahan lan perkalian ing
$\Domain{M_1}$ lan $\Domain{M_2}$ bisa ditetepaké
kanthi cara kang ora isomorfik. (Iki uga manut téoréma
misuwur saka Vaught: cacah modèl terhitung saka
téori jangkep ora bisa persis~2.)

\begin{prob}
Tuduhna yèn ing modèl aritmetika nonstandar ora ana
blok paling gedhé.
\end{prob}

\begin{prob} 
Anggep $\Lang{L}$ sawijining !!{language} tataran kapisan
kang mung ngemot $<$ minangka !!{predicate}
(saliyané $\eq$), lan anggep $\Struct{N} = (\Nat, <)$.
Kabèh himpunan pérangan winates utawa kang komplemèné
winates ing $\Struct{N}$ bisa ditetepaké.
Tuduhna yèn \emph{mung} himpunan mau kang bisa ditetepaké
ing $\Struct{N}$.
 
(Pituduh: Dhisik, anggep $\Obj{prc}(x,y)$ minangka
cekakan formula $\Lang{L}$ kang ateges ``$x$ iku
pendahulu langsung saka $y$:''
\[
x<y \land \lnot \lexists[z][(x<z \land z < y)].
\]
Saiki, saben himpunan pérangan $\Struct{N}$ kang bisa
ditetepaké cocog karo formula $!A(x)$ ing $\Lang{L}$.
Kanggo saben $!A$ kaya mangkono, delengen ukara $!D$:
\[
\lexists[x][\lforall[y][\lforall[z][((x<y \land x<z \land \Obj{prc}(y,z)
  \land !A(y)) \lif !A(z))]]].
\]
Tuduhna yèn $\Sat{N}{!D}$ yèn lan mung yèn
himpunan pérangan $\Struct{N}$ kang ditetepaké déning
$!A$ iku winates utawa komplemèné winates.

Saiki, anggep $\Struct{M}$ modèl nonstandar kang
ekuivalen sacara elementer karo $\Struct{N}$.
Yèn $a \in \Domain{M}$ nonstandar, pilih
$b, c \in \Domain{M}$ kang luwih gedhé tinimbang $a$,
lan anggep $b$ pendahulu langsung saka~$c$.
Mula ana automorfisme $h$ saka $\Domain{M}$
saéngga $h(b)=c$ (kenapa?).
Dadi, yèn $b$ nyukupi $!A$, $c$ uga mangkono
(kenapa?). Akibaté, $!D$ bener ing
$\Struct{M}$ lan mula uga ing $\Struct{N}$.
Iki nuduhaké yèn himpunan pérangan $\Struct{N}$
kang ditetepaké déning $!A$ iku winates utawa
komplemèné winates.
\end{prob}

\end{document}
